\section{问题3：多阶段预报更新与付费调整模型}

\subsection{机制与费用函数}

问题3 中每天 0:00、6:00、12:00、18:00 可获得未来 24 小时整点光伏预报。以 0:00 预报制定全天计划 $\hat g_t$，并在其余三个时刻依据新预报调整当日剩余时段的购电量 $\tilde g_t$；储能计划可随调整同步修改（题目只对购电量偏差计费）。费用规则为：调整量低于计划量的部分按交易时刻电价的 $0.5$ 倍计违约费，高于计划量的部分按 $1.5$ 倍计价，仍不足覆盖实际负载的部分按 $5$ 倍紧急购电。

记 $\Delta_t=\tilde g_t-\hat g_t$，单时段费用为

\begin{equation}
  \Phi_t(\tilde g_t;\hat g_t)=p_t\hat g_t+
  \begin{cases}
    1.5\,p_t\Delta_t, & \Delta_t\ge0,\\[2pt]
    0.5\,p_t\Delta_t, & \Delta_t<0 .
  \end{cases}\text{。}
  \label{eq:p3-phi}
\end{equation}

\textbf{命题1}：式 \eqref{eq:p3-phi} 等价于对称罚金形式

\begin{equation}
  \Phi_t(\tilde g_t;\hat g_t)=p_t\tilde g_t+0.5\,p_t\big|\tilde g_t-\hat g_t\big| .
  \label{eq:p3-sym}
\end{equation}

证明：当 $\Delta_t\ge0$ 时，右端为 $p_t(\hat g_t+\Delta_t)+0.5p_t\Delta_t=p_t\hat g_t+1.5p_t\Delta_t$；当 $\Delta_t<0$ 时，右端为 $p_t(\hat g_t+\Delta_t)+0.5p_t(-\Delta_t)=p_t\hat g_t+0.5p_t\Delta_t$，二者分别与 \cref{eq:p3-phi} 的两支一致。\hfill$\square$

该化简揭示了三档价格机制的本质：\textbf{按实际购电量付 $p_t$，再对偏离计划的部分收取对称的 $0.5p_t$ 罚金}。由此可立即得到两个性质：其一，罚金斜率两侧相同，不存在"故意多报计划再取消"的套利空间（多报 1 \si{kWh} 再取消需付 $0.5p_t$）；其二，目标函数关于 $\tilde g_t$ 是凸的分段线性函数，故整个多阶段问题仍可写成线性规划，无需整数变量。

三档价格的决策含义为：调低 1 \si{kWh} 相对"按计划买"省 $0.5p_t$，调高 1 \si{kWh} 多付 $0.5p_t$，而紧急购电 1 \si{kWh} 多付 $4p_t$。因此\textbf{永远优先"调高"而不是"紧急"}$(1.5p_t\ll5p_t)$，紧急购电只应出现在最后一次预报之后仍然发生的光伏与负载偏差上。

\subsection{滚动时域重优化模型}

\textbf{阶段一（0:00 计划）}：以 0:00 预报 $\hat P^{0}_t$ 与负载预测 $\hat L_t$（叠加问题2 的安全裕度 $\alpha$）解 \cref{eq:p2-sub}，得计划 $(\hat g_t,\hat a_t,\hat b_t)$ 与储电量轨迹；该滚动重优化框架属模型预测控制的思想\cite{mayne2000}。

\textbf{阶段二（$6{:}00/12{:}00/18{:}00$ 调整）}：设调整时刻对应首时段 $t_0=6u+1$（$u=6,12,18$），已知当前实际储电量 $E_{t_0-1}$、承诺计划 $\hat g_t$ 与最新预报 $\hat P^{u}_t$，解

\begin{equation}
  \min_{\{\tilde g_t,a_t,b_t,s_t,E_t\}_{t\ge t_0}}
  \sum_{t=t_0}^{T}\Big[p_t\tilde g_t+0.5\,p_t\big|\tilde g_t-\hat g_t\big|\Big]\text{。}
  \label{eq:p3-stage}
\end{equation}

\begin{equation}
  \text{s.t.}\quad
  \begin{cases}
    \tilde g_t+\eta_db_t+\hat P^{u}_t\Delta t=L_t\Delta t+\dfrac{a_t}{\eta_c}+s_t,&\quad t\ge t_0,\\[4pt]
    E_t=E_{t-1}+a_t-b_t,\quad1200\le E_t\le10800,\quad E_{144}=E_0,&\quad t\ge t_0,\\[4pt]
    \tilde g_t\ge0,\quad0\le a_t,b_t\le833.33,\quad s_t\ge0,&\quad t\ge t_0 .
  \end{cases}\text{。}
  \label{eq:p3-cons}
\end{equation}

式 \eqref{eq:p3-stage} 中的绝对值项用辅助变量精确线性化：令 $\tilde g_t-\hat g_t=\delta_t^{+}-\delta_t^{-}$、$\delta_t^{\pm}\ge0$，以 $\delta_t^{+}+\delta_t^{-}$ 替代 $|\tilde g_t-\hat g_t|$。由于两项系数均为正，最小化会自动使 $\min(\delta_t^{+},\delta_t^{-})=0$，故与 \eqref{eq:p3-phi} 完全等价、不引入松弛误差。已执行时段（$t<t_0$）的计划与实绩被冻结，只有剩余时段参与优化；各阶段的终端条件 $E_{144}=E_0$ 保持一致，避免调整导致跨日漂移。

\subsection{理论分析}

\textbf{命题2（最优计划分位点为中位数）}：设某时段最终所需的购电量 $D$ 为随机量、0:00 计划为 $q$。若在调整时刻能够（近似）获知真实需要量并取 $\tilde g=D$，则该时段费用为 $C(q,D)=pD+0.5p|D-q|$，于是

\begin{equation}
  \min_q\ \mathbb{E}[C]=p\,\mathbb{E}[D]+0.5p\min_q\mathbb{E}|D-q|
  =p\,\mathbb{E}[D]+0.5p\,\mathrm{MAD}(D),
  \label{eq:p3-median}
\end{equation}

其中 $\mathrm{MAD}(D)=\mathbb{E}|D-\mathrm{median}(D)|$ 为平均绝对偏差，最小值在中位数处取得。\textbf{即：调整机制下的最优计划分位点是中位数}。

这一结论与问题2 形成鲜明对比：问题2 中缺电代价为 $4p_t$ 而多备代价仅 $0.19p_t\sim p_t$，最优计划覆盖 $80\%$--$96\%$ 分位（偏保守）；问题3 中偏离计划的惩罚关于计划量对称（两侧均为 $0.5p_t$），故最优计划取中位数（中性）。\textbf{调整机制的存在改变了计划的分位点选择}，这正是"是否需要保守计划"这一问题的答案。

\textbf{命题3（更新价值的分解）}：设某时段在更新前后的不确定性分别为 $D^{\tau-}$ 与 $D^{\tau}$，由 \cref{eq:p3-median} 得

\begin{equation}
  V_\tau=0.5\,p\Big[\mathrm{MAD}(D^{\tau-})-\mathrm{MAD}(D^{\tau})\Big]
  +\Delta\{\text{紧急购电费}\},
  \label{eq:p3-value}
\end{equation}

即更新的价值约等于"$0.5\times$ 电价 $\times$ 平均绝对偏差的下降量"加上紧急购电费的下降量。该式为下文的定量判据提供了理论依据：只有当预报更新能显著降低偏差（或显著减少紧急购电）时，更新才值得做。

\subsection{预报精度与修订幅度}

按\S\ref{sec:boundary} 的口径（预报 $k$ 小时为发布时刻之后的第 $k$ 小时）核算预报精度，并在\textbf{同一时间窗}上比较不同发布时刻的精度，结果见 \cref{tab:fc-acc}。

\begin{table}[H]
  \centering
  \caption{不同发布时刻的预报精度（同一时间窗比较，单位：\si{kW}）}
  \label{tab:fc-acc}
  \small
  \begin{tabular}{lccc}
    \toprule[1.5pt]
    时间窗 & 发布时刻 & MAE & 平均偏差 \\
    \midrule[1pt]
    \multirow{2}{*}{07:00--24:00} & 0:00 & 442.2 & $-47.9$ \\
                                  & 6:00 & 388.0 & $-23.4$ \\
    \midrule[1pt]
    \multirow{3}{*}{13:00--24:00} & 0:00 & 350.5 & $-293.3$ \\
                                  & 6:00 & 304.9 & $-274.2$ \\
                                  & 12:00 & 272.0 & $-262.6$ \\
    \midrule[1pt]
    \multirow{4}{*}{19:00--24:00} & 0:00 & 17.8 & $-17.8$ \\
                                  & 6:00 & 17.8 & $-17.8$ \\
                                  & 12:00 & 17.8 & $-17.8$ \\
                                  & 18:00 & 17.8 & $-17.8$ \\
    \bottomrule[1.5pt]
  \end{tabular}
\end{table}

\cref{tab:fc-acc} 给出三点关键信息：（1）\textbf{后续发布时刻确实更准确}——13:00--24:00 窗口上 MAE 由 0:00 的 $350.5$ \si{kW} 降到 12:00 的 $272.0$ \si{kW}；（2）预报存在系统性偏低，午后低估约 $260$--$293$ \si{kW}；（3）\textbf{夜间 19:00--24:00 四个发布时刻的预报完全相同}（MAE 均为 $17.8$ \si{kW}，均等于平均偏差），因为该时段光伏出力为零，任何预报都无法提供新信息。第三点直接解释了后文 18:00 更新价值为零的实测结果。

相应地，相邻两次预报对同一时钟小时的修订幅度为：0:00$\to$6:00 的 MAE 为 $279.5$ \si{kW}、6:00$\to$12:00 为 $120.0$ \si{kW}、12:00$\to$18:00 仅 $0.002$ \si{kW}（对应夜间时段，实质上无修订）。\cref{fig:p3-rev}(a) 给出了三个修订幅度的对比。

\subsection{求解结果与策略对比}

安全裕度按 1 月数据在 P3 策略下标定为 $\alpha=0.05$（其余策略共用同一裕度，以保证比较公平），报告期结果见 \cref{tab:p3-pol} 与 \cref{fig:p3-pol}。

\begin{table}[H]
  \centering
  \caption{问题3 各策略的全年费用对比（334 天，单位：元）}
  \label{tab:p3-pol}
  \small
  \begin{tabular}{llrrrr}
    \toprule[1.5pt]
    策略 & 说明 & 计划购电费 & 违约费 & 紧急购电费 & 总费用 \\
    \midrule[1pt]
    P0 & 仅 0:00 计划            & 13700904 & 0       & 7551323 & 21252227 \\
    P1 & 0:00$+$6:00 调整         & 13651925 & 327619  & 7442454 & 21421999 \\
    P2 & 0:00$+$6:00$+$12:00      & 13637359 & 303946  & 7428169 & 21369474 \\
    \textbf{P3} & \textbf{0:00$+$6/12/18（题目方案）} & 13637358 & 303947 & 7428170 & \textbf{21369475} \\
    P4 & 逐小时更新（插值预报）   & 13603974 & 303663  & 7477610 & 21385247 \\
    \midrule[1pt]
    P5$'$ & 仅完美光伏信息（单阶段） & 13399959 & --- & 5849762 & 19249721 \\
    A\_load & 仅完美负载信息（单阶段） & 12556455 & --- & 3491877 & 16048332 \\
    P5 & 完美信息（负载$+$光伏）  & 12229643 & --- & 0 & 12229643 \\
    \bottomrule[1.5pt]
  \end{tabular}
  \par\vspace{3pt}{\footnotesize 注：P5$\prime$、A\_load、P5 为单阶段完美信息边界，不涉及计划调整，故不计违约费，其"总费用"等于计划购电费与紧急购电费之和。}
\end{table}

\begin{figure}[H]
  \centering
  \includegraphics[width=0.66\linewidth]{fig11.png}
  \caption{问题3 的策略对比：(a) 各预报更新策略的全年总费用；(b) 计划费、违约费与紧急购电费的构成}
  \label{fig:p3-pol}
\end{figure}

实测结果给出了一个与直觉相反、但对决策极为重要的结论：\textbf{预报更新并未降低成本}。P0 的总费用为 $21252227$ 元，而题目给定的 P3 方案为 $21369475$ 元，反而高出 $117247$ 元（$+0.55\%$）；逐小时更新（P4）亦无增益（较 P3 高 $15772$ 元）；P3 全年紧急购电量为 $169.97$ 万 \si{kWh}，全部来自最后一次预报之后仍未消除的残差。其机理可由 \cref{eq:p3-value} 分解：6:00 更新带来的紧急购电费下降约 $10.9$ 万元，但按 $0.5p_t$ 计收的违约费达到 $32.8$ 万元，净效应为负；12:00 更新的边际价值为正（$+5.25$ 万元，因为该时段预报修订 $120.0$ \si{kW} 且覆盖午后光伏出力时段），18:00 更新的边际价值为零（$-0.76$ 元，因为夜间预报无修订）。各更新时点的边际价值见 \cref{fig:p3-rev}(b)。

\begin{figure}[H]
  \centering
  \includegraphics[width=0.66\linewidth]{fig12.png}
  \caption{问题3：(a) 相邻预报时刻的修订幅度；(b) 信息价值分解（单阶段口径）}
  \label{fig:p3-rev}
\end{figure}

\subsection{是否需要引入其他时刻的预报}

题目要求分析是否需要在 0:00、6:00、12:00、18:00 之外引入其他时刻的预报来制定调整策略。本文从三个层次回答。

\textbf{（1）边际价值角度。} 各更新时点的实测边际价值为：6:00 为 $-16.98$ 万元，12:00 为 $+5.25$ 万元，18:00 为 $-0.08$ 万元；逐小时更新相对题目方案为 $-1.58$ 万元。即除 12:00 外，增加更新均不产生正收益；而逐小时更新（信息粒度进一步细化）也不再改善。这说明\textbf{在本题数据条件下 4 个时点已接近饱和}。

\textbf{（2）信息量角度。} 由 \cref{tab:fc-acc}，夜间时段（19:00--24:00）四个发布时刻的预报完全相同，18:00 更新不含任何新信息，其价值必然为零；而白天时段的预报修订幅度也仅 $120$--$280$ \si{kW}，相对负载预测误差（7 日滑动平均 MAE $781.2$ \si{kW}、RMSE $955.9$ \si{kW}）小得多。因此即使加密预报时刻，可改善的也只是"光伏侧的小误差"，无法触及主要误差来源。

\textbf{（3）信息价值角度。} 用完美信息边界分解两类信息价值：仅完美光伏信息时总费用为 $19249721$ 元，价值 $200.25$ 万元；仅完美负载信息时总费用为 $16048332$ 元，价值 $520.39$ 万元（\cref{fig:p3-rev}(b)）；若负载与光伏均完美已知，总费用可降至 $12229643$ 元，相对 P0 的信息价值为 $902.26$ 万元。\textbf{负载信息的价值是光伏信息的 2.6 倍}，而题目并未提供负载预报产品，其误差在本模型框架内不可改善。

\textbf{结论}：\textbf{在本题数据条件下，不需要引入其他时刻的预报来制定调整购电策略}。预报更新机制本身并未降低总费用，其原因是（i）可更新的光伏信息价值有限，（ii）调整需按 $0.5p_t$ 支付罚金，（iii）主要不确定性来自负载侧。若要进一步提升经济效益，应优先获取\textbf{负载预报}（如基于气象与用电行为的中短期负荷预测），或推动制度层面降低调整费率；若必须在本模型内使用更新信息，则建议只保留 12:00 时点，因其边际价值为正。

需要说明的是，上述结论依赖于题目给定的预报数据质量。本文将附件3 预报误差人为放大 $1.5$ 倍后重做全部策略（见 \cref{tab:sens-fc}）：误差放大使各策略费用普遍上升，其中 12:00 更新的边际价值由 $5.25$ 万元升至 $303.82$ 万元，说明	extbf{预报质量越差、及时更新越有价值}；但 6:00 更新需支付的违约费同步放大，P3 相对 P0 的差额由 $+0.55\%$ 扩大到 $+9.97\%$。因此"更新无正收益"的结论在预报质量恶化时依然成立，且在数据质量较好时更为明确。
